\bookchapter{Capitolo 27: Il telefono facile anche per i genitori — Guida completa all'adattamento per gli anziani}

\chapterintro{A dire il vero, mentre scrivevo questo capitolo, ho pensato ai miei genitori. Ogni volta che sul telefono di mia madre compare “aggiornamento di sistema”, non osa cliccare, per paura di romperlo. Mio padre ha il carattere di WeChat più piccolo delle formiche, ma non sa come ingrandirlo, e così ha strizzato gli occhi per due anni. Non sono “stupidi”, è che i telefoni, per default, trattano tutti come ventenni. In questo capitolo, con parole semplici e dirette, sciogliamo uno per uno i problemi più fastidiosi che gli anziani incontrano nell'uso del telefono.}

\booksection{Prima di tutto, l'atteggiamento giusto: non è che “gli anziani non sanno usare”, è che “il telefono non è ben fatto”}

Prima di qualsiasi impostazione, c'è un'idea da chiarire:\textbf{le interfacce degli smartphone di oggi, fin dalla progettazione, hanno considerato poco gli utenti anziani}. Pop-up piccoli che non si leggono, gesti che richiedono tocchi precisi, icone tutte uguali — questi non sono problemi degli utenti, sono problemi di prodotto. Quindi, aiutare gli anziani con il telefono, la prima cosa non è insegnare, ma aiutarli\textbf{a configurare bene l'ambiente}. Una volta configurato l'ambiente, la maggior parte dei problemi scompare.

Ci sono tre regole d'oro per aiutare gli anziani con il telefono: (1)\textbf{Non avere fretta}— tu usi smartphone da dieci anni, loro magari da tre mesi. (2)\textbf{Usa analogie}— “il Wi-Fi è come l'acqua di casa, apri il rubinetto e l'acqua scorre”, queste metafore sono molto più comprensibili del gergo tecnico. (3)\textbf{Scrivi}— a voce se lo dimenticano, ma se scrivi i passaggi chiave su un foglio e lo attacchi dietro la cover del telefono, funziona meglio di qualsiasi altra cosa.

\booksection{Primo passo: attivare la “Modalità semplice” — il telefono diventa pulito}

I telefoni di oggi integrano quasi tutti una “Modalità semplice” (chiamata da alcuni “Modalità assistenza”, “Modalità anziani”, “Modalità anziani”), una volta attivata, l'intero telefono cambia aspetto:

\begin{itemize}
  \item \textbf{Caratteri e icone diventano grandi}— non serve più la lente d'ingrandimento.
  \item \textbf{Le icone del desktop sono molto semplificate}— restano solo le App più usate come Telefono, Messaggi, Fotocamera, Galleria, WeChat; tutto il resto viene nascosto.
  \item \textbf{Le operazioni si semplificano}— molte funzioni vistose (schermata -1, assistente intelligente, finestre fluttuanti) vengono chiuse, l'interfaccia punta alla pulizia.
  \item \textbf{L'audio aumenta}— suoneria, volume notifiche, audio delle chiamate vengono regolati automaticamente al massimo.
\end{itemize}

\booksubsection{Metodi di attivazione della modalità semplice per brand (segui direttamente le istruzioni):}

\begin{center}
\begin{tabularx}{\linewidth}{|X|X|X|}
\hline
\textbf{Brand} & \textbf{Nome modalità} & \textbf{Come attivarla} \tabularnewline
\hline
iPhone & Accessibilità (nessuna modalità semplice separata, ma può essere ingrandita voce per voce) & Impostazioni > Display e luminosità > Dimensione testo, trascina al massimo. Quindi in "Accessibilità > Display e dimensione testo" attiva "Testo in grassetto" e "Forma dei pulsanti" \tabularnewline
Huawei/Honor & Modalità semplice & Impostazioni > Sistema e aggiornamenti > Modalità semplice > Attiva \tabularnewline
Xiaomi/Redmi & Modalità essenziale & Impostazioni > Altre impostazioni > Modalità essenziale > Attiva \tabularnewline
OPPO/OnePlus & Modalità semplice & Impostazioni > Funzioni speciali > Modalità semplice \tabularnewline
vivo/iQOO & Desktop semplice & Impostazioni > Desktop, schermata di blocco e sfondo > Desktop semplice \tabularnewline
Samsung & Modalità semplice & Impostazioni > Display > Modalità semplice \tabularnewline
\hline
\end{tabularx}
\end{center}

Se usi un altro brand, cerca “semplice” o “anziani” nella casella di ricerca in alto nelle Impostazioni del telefono: di solito lo trovi. Se non lo trovi, segui la soluzione “ingrandimento manuale” qui sotto, voce per voce: l'effetto è lo stesso.

\booksection{Se non attivi la modalità semplice, anche il manuale funziona — “Caratteri grandi + icone grandi” in cinque passi}

Alcuni anziani non vogliono usare la modalità semplice, la trovano troppo “vecchia”. Allora regola manualmente queste cinque voci, con lo stesso risultato:

\begin{enumerate}
  \item \textbf{Caratteri al massimo}: Impostazioni > Display > Dimensione caratteri, trascina tutto a destra. Approfitta per mettere anche "Dimensione display" (a volte "Scala display") al massimo, così tutti gli elementi dell'interfaccia si ingrandiscono.
  \item \textbf{Icone più grandi}: nella schermata Home, stringi due dita (come per rimpicciolire una foto) per entrare nelle Impostazioni desktop > Layout desktop (o Dimensione icone), scegli il layout con meno celle (ad esempio 3×5), le icone saranno naturalmente più grandi.
  \item \textbf{Attiva "Testo ad alto contrasto"}: rende i bordi del testo più nitidi; lo trovi in "Accessibilità".
  \item \textbf{Disattiva le notifiche non necessarie}: Impostazioni > Notifiche, disattiva tutto tranne Telefono, Messaggi e WeChat. Così il telefono non suona ogni momento e non infastidisce.
  \item \textbf{Metti i contatti frequenti sulla schermata Home}: premi a lungo su un punto vuoto del desktop > Aggiungi widget > Contatti, metti sulla schermata Home i contatti rapidi per coniuge, figli, medico, contatti di emergenza. Basta un tocco per chiamare, senza scorrere la rubrica.
\end{enumerate}

\begin{tcbpractice}{Fai: aiuta gli anziani a "ristrutturare" il telefono}
\begin{enumerate}
  \item Attiva la modalità semplice (segui la tabella dei brand sopra).
  \item Attiva la visualizzazione numerica della batteria (Impostazioni > Batteria > Percentuale batteria), così sanno sempre quanta carica resta.
  \item Apri il pannello notifiche e rimuovi o nascondi i toggle rapidi che gli anziani non usano, come Bluetooth, NFC, modalità aereo.
  \item Metti quattro contatti rapidi sul desktop: coniuge, figli, medico di famiglia (se c'è), community/condominio.
\end{enumerate}

\end{tcbpractice}

\booksection{Prevenzione delle truffe — leggi questo capitolo insieme agli anziani}

Gli anziani sono la categoria più colpita dalle truffe telecom. Non è che abbiano poca vigilanza, è che i copioni dei truffatori sono studiati con cura, mirati alle debolezze psicologiche degli anziani. Ecco i tipi di truffa più comuni e come affrontarli: ti consiglio di leggerli ad alta voce agli anziani o di stamparli e attaccarli vicino al telefono.

\booksubsection{Sei truffe comuni (se le vedi, riaggancia subito):}

\begin{itemize}
  \item \textbf{1. "Tuo nipote ha avuto un problema, manda soldi subito"}— il truffatore dice "tuo nipote/nipotina ha combinato un guaio fuori, è stato arrestato dalla polizia, serve una cauzione". Ricorda:\textbf{qualsiasi chiamata che richiede un bonifico per "far rilasciare qualcuno" è una truffa}. La polizia e la magistratura non chiedono soldi per telefono.
  \item \textbf{2. "La tua carta sanitaria è stata sospesa"}— un SMS dice che la carta è anomala e ti fa cliccare su un link per "attivarla". Ricorda:\textbf{l'ente sanitario e l'INPS non mandano SMS con link da cliccare}. In questi casi chiama direttamente il numero ufficiale per verificare.
  \item \textbf{3. "Rendimento alto, guadagno sicuro"}— nei gruppi WeChat o via SMS propongono "investimenti pensionistici" o "progetti statali", promettendo rendimenti elevati. Ricorda:\textbf{chi promette rendimenti annui sopra il 5\% con "capitale garantito", nel 99,9\% dei casi è una truffa}.
  \item \textbf{4. "Sono il tuo capo/vecchio amico, ho cambiato numero"}— si spacciano per conoscenti e chiedono soldi in prestito. Ricorda:\textbf{una telefonata di verifica vale più di qualsiasi cosa}, chi dice "non posso rispondere al telefono" per l'80\% è un truffatore.
  \item \textbf{5. "Uova gratis / conferenze su prodotti sanitari"}— eventi gratuiti offline attirano gli anziani a una "conferenza salute" e poi vendono integratori a prezzi esorbitanti. Ricorda:\textbf{non cade pere dal cielo, dietro le iniziative gratuite c'è sempre un business}.
  \item \textbf{6. "Il tuo telefono è infetto, serve una riparazione remota"}— il truffatore ti fa scaricare un'App di controllo remoto (come Sunflower, TeamViewer), poi prende il controllo del telefono e trasferisce i soldi. Ricorda:\textbf{non permettere mai a nessuno di controllare il tuo telefono da remoto!}I produttori di telefoni, gli operatori telefonici e le banche non chiamano per farti installare software di controllo remoto.
\end{itemize}

\booksubsection{Tre linee di difesa: una volta attivate, bloccano il 90\% delle truffe:}

\begin{enumerate}
  \item \textbf{Prima linea: blocco nativo del telefono}— sui telefoni Android attiva il "blocco chiamate fastidiose" (nell'App Telefono o in Phone Manager); su iPhone attiva "Silenzia numeri sconosciuti" (Impostazioni > Telefono > Silenzia numeri sconosciuti).
  \item \textbf{Seconda linea: blocco dell'operatore}— invia l'SMS "KTFSR" al 10086 per attivare gratuitamente il blocco delle chiamate fastidiose ad alta frequenza di China Mobile; anche Telecom e Unicom hanno servizi simili. Chiama il servizio clienti dell'operatore e chiedi l'attivazione.
  \item \textbf{Terza linea: App del Centro nazionale anti-truffa}— scarica dallo store l'App "Centro nazionale anti-truffa", registrati e attiva l'avviso di chiamate. Questa App è collegata al database della polizia, riconosce i numeri dei truffatori non appena chiamano.
\end{enumerate}

\begin{tcbwarn}{La "filastrocca anti-truffa" per gli anziani (stampala e attaccala vicino al telefono)}
\begin{itemize}
  \item \textbf{Chiamata sconosciuta che chiede soldi riaggancia!}
  \item \textbf{SMS con link da cliccare cancellalo!}
  \item \textbf{Promesse di rendimenti alti con capitale garantito blocca!}
  \item \textbf{Ti fa installare software di controllo remoto chiudi subito!}
  \item \textbf{La vera polizia non indaga per telefono e non chiede soldi!}
  \item \textbf{In caso di dubbio, prima chiama i figli e chiedi.}
\end{itemize}

\end{tcbwarn}

\booksection{Assistenza remota — anche da lontano puoi sistemare il telefono dei genitori}

A volte gli anziani incontrano problemi che al telefono non riesci a descrivere: non trovano il pulsante. Qui entra in gioco l'“assistenza remota” — puoi vedere lo schermo del loro telefono direttamente sul tuo e aiutarli a operare. Le soluzioni variano per brand:

\begin{itemize}
  \item \textbf{Huawei/Honor}: l'App nativa "MeeTime" ha una funzione "Condivisione schermo". Se tu e i tuoi genitori avete entrambi MeeTime, dopo una videochiamata cliccate su "Condividi schermo", puoi disegnare cerchi sul loro schermo e dire a tua madre "clicca qui".
  \item \textbf{Xiaomi/Redmi}: nell'App Mi Video Call c'è "Assistenza remota", con condivisione schermo + chiamata vocale sincronizzata.
  \item \textbf{OPPO/vivo}: integra "Family Space" o "Remote Guard", cercando "remoto" nelle Impostazioni lo trovi.
  \item \textbf{Soluzione universale (per tutti i brand)}: scaricate l'App "Sunflower Remote Control" (la versione gratuita basta), installatela entrambi. Con lo stesso account o aggiungendovi come amici, fate leggere ai genitori il codice di verifica mostrato sul loro telefono, e potrai vedere e operare a distanza il loro schermo.
  \item \textbf{Utenti iPhone}: FaceTime supporta la "condivisione schermo" durante la chiamata. In FaceTime, clicca sul pulsante di condivisione schermo e l'altro vedrà il tuo schermo. Ma per limiti di sicurezza di iOS non si può operare a distanza, solo "insegnamento per indicazione".
\end{itemize}

\booksection{Le “modifiche per anziani” di WeChat — bastano queste frasi da dire ai genitori}

WeChat è l'App più usata dagli anziani, ma l'interfaccia è troppo fitta e complessa. Ecco alcune modifiche con effetto immediato:

\begin{itemize}
  \item \textbf{Ingrandisci i caratteri interni di WeChat}: WeChat > Me > Impostazioni > Generali > Dimensione testo, trascina al massimo. Nota: influisce solo sui caratteri interni di WeChat, non su quelli di sistema.
  \item \textbf{Attiva la "Modalità assistenza"}: WeChat > Me > Impostazioni > Modalità assistenza > Attiva. Una volta attiva, caratteri e pulsanti si ingrandiscono, aumenta il contrasto dei colori, e si può "ascoltare il messaggio di testo" con un tocco (cliccando su un messaggio di testo viene letto ad alta voce).
  \item \textbf{Personalizza i contatti importanti}: cambia in rubrica il nome dei figli e dei parenti in modo riconoscibile — ad esempio "Figlio grande - Ming", "Seconda figlia - Hong", così aprono WeChat e sanno subito chi è chi.
  \item \textbf{Insegna loro a usare il "dettatura vocale" invece di digitare}: sulla tastiera accanto alla chat, tieni premuta la barra spaziatrice (Android) o tocca l'icona del microfono (iPhone), parla e il testo viene inviato. La precisione è ormai altissima, per gli anziani che non conoscono il pinyin è una comodità "rivoluzionaria".
  \item \textbf{Disattiva le funzioni inutilizzate in "Scopri"}: WeChat > Me > Impostazioni > Generali > Gestione pagina Scopri, disattiva tra "Momenti", "Canali video", "Live", "Shopping", "Giochi" quelle che non usi, per ridurre le distrazioni.
\end{itemize}

\booksection{Cosa fare se si perde il telefono? Prepara con gli anziani un “piano di emergenza”}

Prevenire è meglio che curare. Finché il telefono non è perso, fai queste cose in anticipo:

\begin{enumerate}
  \item \textbf{Verifica che "Trova il mio telefono" sia attivo}: iPhone in Impostazioni > [tuo nome] > Trova > Trova il mio iPhone; Android in Impostazioni cerca "Trova dispositivo". Assicurati che sia attivo.
  \item \textbf{Segna le password e gli account importanti su un foglio}: la password di Apple ID / account Huawei / account Xiaomi, una copia chiusa in un cassetto, una copia a un figlio fidato. Se si perde il telefono, serve accedere a questi account per il blocco e la localizzazione remoti; se si dimentica la password, sono guai.
  \item \textbf{Insegna agli anziani una frase}:\textbf{"Se perdi il telefono, per prima cosa chiama un familiare."}Il familiare può accedere subito a iCloud o al sito Trova dispositivo di Android, contrassegnarlo remotamente come "Modalità smarrito", così chi lo trova non può aprirlo né usarlo.
\end{enumerate}

\booksection{Compatibilità auricolari wireless e apparecchi acustici — una funzione sottovalutata}

Molti anziani hanno problemi di udito, ma non sanno che il telefono può funzionare insieme all'apparecchio acustico. iPhone e alcuni Android supportano gli apparecchi acustici MFi (Made for iPhone), possono inviare direttamente l'audio della chiamata all'apparecchio, molto più chiaro che tenere il telefono all'orecchio.

\begin{itemize}
  \item \textbf{Metodo di collegamento}: iPhone Impostazioni > Accessibilità > Apparecchi acustici, attiva il Bluetooth e segui le istruzioni di abbinamento dell'apparecchio acustico.
  \item \textbf{Se non usi l'apparecchio acustico}: in "Accessibilità > Audio/Visivo" attiva "Riduzione rumore telefono" e "Audio mono", migliora la chiarezza della chiamata.
  \item \textbf{Android}: in "Accessibilità > Potenziamento uditivo" ci sono impostazioni simili, la posizione varia leggermente per brand; cerca direttamente "apparecchio acustico".
\end{itemize}

\booksection{Funzioni legate alla salute — il telefono non è solo uno strumento di comunicazione}

\begin{itemize}
  \item \textbf{SOS di emergenza}: premi il pulsante di accensione cinque volte di fila (sulla maggior parte dei telefoni), il telefono chiama automaticamente il numero di emergenza e invia la posizione ai contatti di emergenza preimpostati. Aiuta gli anziani a configurarlo in anticipo — iPhone: "Impostazioni > SOS di emergenza", Android: "Impostazioni > Sicurezza > SOS di emergenza". Fai una prova con gli anziani prima di insegnar loro, così non sbagliano per la tensione del momento.
  \item \textbf{Rilevamento cadute}: Apple Watch ha questa funzione, se rileva una caduta e resti fermo per un po', chiama automaticamente i soccorsi e avvisa la famiglia. Se il budget lo permette, regalalo agli anziani: nei momenti cruciali può salvare la vita.
  \item \textbf{Promemoria per i farmaci}: l'App "Salute" di iPhone ha una funzione di promemoria farmaci integrata, impostando nome del farmaco, orario e dose, alla scadenza compare un avviso. Su Android puoi usare "allarmi programmati" o scaricare un'App tipo "promemoria farmaci".
\end{itemize}

\begin{tcbsummary}{Riepilogo del capitolo}
\begin{itemize}
  \item Prima attiva la modalità semplice o ingrandisci manualmente caratteri/icone: ambienti ben configurati risolvono la maggior parte dei problemi
  \item Tre regole d'oro per aiutare gli anziani con il telefono: non avere fretta, usa analogie, scrivi
  \item Sei truffe comuni da ricordare sempre; tre linee di difesa bloccano il 90\% delle truffe
  \item L'assistenza remota ti permette di sistemare il telefono dei genitori anche da lontano
  \item WeChat con "Modalità assistenza" + "dettatura vocale" sono le due funzioni più pratiche per gli anziani
  \item SOS di emergenza e rilevamento cadute: configurazioni che salvano la vita, valgono cinque minuti di setup
\end{itemize}

\end{tcbsummary}
